\documentclass[11pt,letterpaper]{article}
% Shared declarations; root documents always import a concrete semantic module.
\usepackage{business-admin}
\usepackage{amsmath,amssymb}
\usepackage{longtable,tabularx,array,booktabs}
\usepackage{enumitem}
\usepackage{ragged2e}
\newcolumntype{Y}{>{\RaggedRight\arraybackslash}X}
\newcommand{\SuiteMeta}[2]{\def\DocTitle{#1}\def\DocSubtitle{#2}\def\DocVersion{1.0}\def\DocStatus{Educational reference}\def\DocOwner{Small Business Management Suite}\def\DocSummary{Original teaching and practice resource; no real-business approval implied.}}
\newcommand{\SuiteOpen}{\makedoctitle\vspace{-0.7em}\noindent\rule{\linewidth}{0.45pt}\par\medskip}
\newcommand{\SourceNote}[2]{\par\medskip\noindent{\footnotesize\textbf{Reading anchor:} Longenecker, Petty, Palich and Hoy, \emph{Small Business Management: Launching \& Growing Entrepreneurial Ventures}, 20th ed., Ch.~#1, textbook p.~#2. This document is independently authored commentary and exercises, not a reproduction of the chapter.}}
\newcommand{\WritingLines}[1]{\par\smallskip\noindent\rule{\linewidth}{0.35pt}\par\vspace{0.35cm}\noindent\rule{\linewidth}{0.35pt}\par\vspace{0.35cm}\noindent\rule{\linewidth}{0.35pt}\par\medskip}
\newcommand{\FormField}[1]{\par\medskip\noindent\textbf{#1}\par\WritingLines{3}}
\newcommand{\SBMCheckBox}{\(\square\)}
\newcommand{\AssessmentRow}[2]{#1 & #2 & \(\square\) & \(\square\) & \(\square\) & \(\square\) & \(\square\)\\\hline}
\newcommand{\AssessmentHeader}{\noindent\begin{tabular}{|p{0.22\linewidth}|p{0.49\linewidth}|c|c|c|c|c|}\hline\textbf{Criterion} & \textbf{Evidence to inspect} & \textbf{0} & \textbf{1} & \textbf{2} & \textbf{3} & \textbf{4}\\\hline}
\newcommand{\AssessmentEnd}{\end{tabular}\par}
\setlength{\emergencystretch}{3em}

\SuiteMeta{Chapter 23: Managing Risk}{Study guide | entrepreneurship and small business management}
\begin{document}
\SuiteOpen
\section{Learning outcome}
Maintain a risk register whose responses match credible exposure and mitigation evidence.
\section{Prerequisites}
Read chapters 2, 8, 18, 20, 21, 22.\par\textbf{Required baseline:} ability to record assumptions and distinguish observations from opinions.
\section{Chapter concept map}
\begin{itemize}[leftmargin=*,itemsep=3pt]
\item business risk
\item property, personnel and liability risk
\item risk identification and treatment
\item insurance selection
\item business continuity
\end{itemize}

\section{Technical study notes}
\begin{itemize}[leftmargin=*,itemsep=3pt]
\item Business risk includes exposures from property, people, liability and interrupted operations, while strategic uncertainty also deserves separate management attention.
\item A risk register is useful only when ratings have rationale, controls are assigned and residual exposure is reviewed after implementation.
\item Insurance may transfer part of financial loss but cannot substitute for prevention, incident response, contractual clarity and continuity rehearsal.
\end{itemize}

\section{Reading and analysis protocol}
\begin{enumerate}[leftmargin=*,itemsep=5pt]
\item Read the chapter learning objectives, then mark terms and claims that need primary evidence.
\item For each chapter theme, distinguish the entrepreneur decision, information required, parties affected and failure condition.
\item Write a one-paragraph synthesis in your own words; cite the textbook chapter and separate contemporary practice notes.
\item Complete the applied activity and attach the underlying evidence, calculations or interview log.
\end{enumerate}

\section{Worked application and deliverable}
\textbf{Business exercise.} Conduct a tabletop on a supplier outage and a customer-data incident.\par\textbf{Required submission.} Risk response and continuity plan. Include a dated assumption register and an explicit proceed/revise/stop recommendation.\par\textbf{Assurance checks.} Risk owner, impact, likelihood, triggers and recovery tests.
\section{Common reasoning failures}
\begin{itemize}[leftmargin=*,itemsep=3pt]
\item Equating insurance with a complete control system
\item Leaving residual risk unowned
\end{itemize}
\noindent\textbf{Correction:} require counterevidence and identify the responsible decision-maker before approving the next stage.
\section{Discussion prompt}
Which risk should be reduced by redesign before transferring it?\par\medskip\textbf{Evidence expected:} at least two alternatives, one quantitative or independently checkable observation, a trade-off, and a decision threshold.
\section{Exit assessment}
\begin{enumerate}[leftmargin=*,itemsep=5pt]
\item Explain each term from the concept map without copying the textbook.
\item Defend the applied decision using an assumption, evidence and counterargument.
\item Show an objective test or inspection that could prove a key assumption wrong.
\item State what would change the recommendation and who has approval authority.
\end{enumerate}

\SourceNote{23}{627}
\end{document}
